% Build: tectonic PlanBee_Architecture.tex
% Korean font: Apple SD Gothic Neo (macOS); replace with an installed Korean font elsewhere.
\documentclass[11pt,a4paper]{article}
\usepackage[margin=23mm]{geometry}
\usepackage{fontspec}
\setmainfont{Apple SD Gothic Neo}
\setsansfont{Apple SD Gothic Neo}
\setmonofont{Menlo}[Scale=0.85]
\usepackage{amsmath,booktabs,tabularx,array,xcolor,tikz,enumitem,fancyhdr,hyperref}
\usetikzlibrary{arrows.meta,positioning,calc}
\definecolor{ink}{HTML}{18334B}
\definecolor{blue}{HTML}{EAF3FC}
\definecolor{green}{HTML}{EAF5EF}
\definecolor{amber}{HTML}{FFF3DE}
\hypersetup{colorlinks=true,linkcolor=ink,urlcolor=ink}
\XeTeXlinebreaklocale "ko"
\XeTeXlinebreakskip=0pt plus 1pt
\setlength{\parindent}{0pt}
\setlength{\parskip}{6pt}
\setlist{nosep,leftmargin=1.5em}
\pagestyle{fancy}\fancyhf{}\fancyhead[L]{\small Plan Bee / V2 Architecture}\fancyhead[R]{\small 2026-11-22}\fancyfoot[C]{\thepage}
\setlength{\headheight}{15pt}
\newcolumntype{Y}{>{\raggedright\arraybackslash}X}
\tikzset{box/.style={draw=ink!60,rounded corners=3pt,fill=blue,align=center,text width=11.5cm,minimum height=9mm,font=\small,inner sep=5pt},arr/.style={-{Latex[length=2mm]},thick,draw=ink}}
\begin{document}
\begin{center}
{\Huge\bfseries Plan Bee}\\[5mm]
{\LARGE 전체 실행 구조와 에이전트 협업 흐름}\\[3mm]
{\large 후보 생성 · 위임카드 기반 토론 · PM 최종 결정}\\[3mm]
2026년 11월 22일 / 현행 V2 구현 기준
\end{center}
\vspace{3mm}
\textbf{핵심 구조.} 사람이 업무 변경 이슈를 줄글로 입력하면, 생성 AI와 알고리즘이 새로운 milestone 계획 후보를 만든다. Hive는 일정과 하드 제약을 코드로 검증한다. 각 개인의 Proxy LLM은 자신의 위임카드로 후보들을 평가하고 공개 토론한다. 추천 후보 1--3개를 받은 PM이 최종 선택하며, 필요한 본인 확인과 최신 Hive 검증이 끝나야 실제 보드에 반영된다.

\textbf{이 문서의 범위.} 현재 서비스의 \texttt{candidate\_discussion} 경로를 설명한다. 이전 \texttt{legacy} 협상 경로는 호환용으로 남아 있으나 본문의 주 흐름과 구분한다. 구현 설명과 실제 Live 성능 검증을 구분하며, 아직 해결되지 않은 응답 오류도 마지막 절에 기록한다. 이 문서 작성에서는 앱 코드 변경이나 Live 실행을 하지 않았다.

\section{주체와 판단 권한}
\begin{tabularx}{\linewidth}{@{}p{27mm}YY@{}}\toprule
주체 & 수행하는 일 & 수행하지 않는 일\\\midrule
사람 사용자 & 역할 선택, 이슈 입력, 본인 확인, PM 역할의 최종 선택 & 채팅 입력만으로 계획 자동 승인\\
생성 AI & Task 내용에 근거한 분해·이동·지원 구조 설계 & 개인 비공개 카드 열람, 공정성 승자 결정\\
탐색 알고리즘 & 담당·시작·분량 조합 확장, 중복 제거 & 최소 부담 분산만으로 적용안 선정\\
Hive / 엔진 & 일정 계산, 제약 검사, 호출 순서·상태·예산 관리 & 개인의 카드상 의사를 대신 결정\\
개인 Proxy LLM & 자기 카드에 따른 평가·발언·거절·동의 & 다른 사람 카드 원문 열람, 마감 제약 무시\\
PM & 추천 후보 선택·전부 반려·수정 요청 & 당사자 거절 자동 무효화, 하드 제약 우회\\\bottomrule
\end{tabularx}

Hive의 일정 검증은 결정적 코드이며 별도의 LLM 심판이 아니다. 자연어 위임 의사는 LLM이 판단하고, 숫자·일정·구조의 유효성은 코드가 검증한다.

\newpage
\section{전체 실행 흐름}
\begin{center}
\begin{tikzpicture}[node distance=4mm]
\node[box] (input) {\textbf{01 사람 입력}\\역할 선택 → 자기 업무 선택 → 줄글 이슈 입력 → 보내기};
\node[box,below=of input] (estimate) {\textbf{담당 Proxy: 소요 변경 여부 판단}\\단순 일정 이동과 실제 작업량 증가 구별 / 초기 보드·카드·프롬프트 고정};
\node[box,below=of estimate] (generate) {\textbf{02 새로운 milestone 후보 생성}\\생성 AI의 의미 설계 + 알고리즘의 담당·분량·시작 조합 확장};
\node[box,fill=green,below=of generate] (validate) {\textbf{Hive: 복제 보드에서 feasibility 검증}\\역할 · 가용 시간 · 용량 · 의존관계 · 마감 · 업무 범위 / 실행 결과 중복 제거};
\node[box,below=of validate] (assess) {\textbf{03 각 개인 Proxy: 카드 기반 평가와 토론}\\개인 평가 → 추천 shortlist → 공개 질의·응답·조건 → 최종 평가};
\node[box,below=of assess] (consent) {\textbf{영향 당사자: 후보별 정확한 버전의 의사 확인}\\ACCEPT / REJECT / NEEDS\_INPUT → 추천 후보 최대 1--3개};
\node[box,fill=amber,below=of consent] (pm) {\textbf{PM 최종 선택 대기}\\하나여도 자동 적용하지 않음 / 비교 후 선택·전부 반려·수정 요청};
\node[box,fill=amber,below=of pm] (human) {\textbf{필요한 당사자 본인 확인}\\각 본인 역할에서 승인·거절 / PM 선택과 개인 동의는 별개};
\node[box,fill=green,below=of human] (commit) {\textbf{Hive 최신 재검증 → 원자적 적용}\\보드·동의·후보 hash 확인 / 실패하면 미합의 변경 미적용};
\foreach \a/\b in {input/estimate,estimate/generate,generate/validate,validate/assess,assess/consent,consent/pm,pm/human,human/commit} \draw[arr] (\a)--(\b);
\end{tikzpicture}
\end{center}
\textbf{그림 1.} 정상 경로의 호출 순서. 공개 조건에 현재 후보로 대응할 수 없으면 제한된 보충 생성과 재평가로 돌아간다. 후보가 없거나 호출·시간 한도에 도달하면 사유를 기록해 PM 검토로 넘기며 적용하지 않는다.

마감이 유효하지 않으면 먼저 기한 회복 후보를 검토하고, 유효한 기준 배분을 중심으로 공정성 논의를 수행한다. 마감이 이미 안전하면 유효한 현재 배분도 유지 후보로 다룰 수 있다. 단계별 snapshot과 적용 여부를 구분해야 하며, 후보 미리보기를 실제 보드로 읽으면 안 된다.

\newpage
\section{시간 안에 가능한 새로운 milestone 설계}
후보는 단순히 ``누가 몇 일을 맡는가''가 아니다. 실제 산출물 범위, 담당자, 작업 순서, 중간 산출물, 의존관계, milestone 시점을 함께 담는 계획이다.

\subsection{생성 AI와 알고리즘의 역할 분리}
생성 AI에는 이슈, Task 설명·완료 조건·분리 경계, 공개 역할과 일정 정보를 전달한다. 개인 비공개 위임카드는 전달하지 않는다. 생성 결과는 구조화된 \texttt{Designs / Transformation}이며, 각 조각은 \texttt{key}, \texttt{what}, \texttt{effort}, \texttt{depends\_on}을 갖는다. 알고리즘은 이 구조를 실제 담당자와 시간 조합으로 확장한다. 명시된 Task 계약을 이용하는 결정적 템플릿도 활용한다.

\begin{tabularx}{\linewidth}{@{}p{28mm}Y@{}}\toprule
변환 & 계획에서 바뀌는 내용\\\midrule
분해·전체 양도 & 의미 있는 하위 산출물을 여러 명이 맡거나 Task 전체를 다른 사람에게 이전\\
일정 이동 & 허용된 미착수 업무의 시작 순서를 바꿔 후속 업무의 시간을 확보\\
협업 & 실제 지원 업무를 배정하고 주 작업 단축의 근거를 남김. Task에서 금지하면 차단\\
후속 이관 & 명시된 \texttt{defer\_option}에 따라 선택 범위를 Beta 등으로 이관. 소요량을 삭제하지 않음\\
선택 작업 생략 & optional 계약, 대체 검증 범위·최소 시간·당사자 확인을 충족할 때만 허용\\
단계 재구성 & 임시 반영·최종 교체처럼 명시된 하위 범위를 만들고 입력 의존성을 해당 단계에 연결\\\bottomrule
\end{tabularx}

분해 시 원 업무 기준 \texttt{effort}와 담당자별 실제 소요일 \texttt{days}를 구분한다. 수치만 나눈 조각은 의미 있는 업무 설계를 보장하지 않는다. 명시 계약은 코드로 검사하되, 범위의 의미상 적절성은 생성 근거와 당사자 토론에도 의존한다.

\subsection{Hive의 검증과 탐색 한계}
모든 후보는 격리된 보드에서 실행해 검사한다. 역할·허용 담당자, 일별 용량·가용 시간, 선행 산출물, Handoff·최종 마감, 업무량 보존, 허용 변경 범위를 검증한다. 사람의 승인이 필요한 0일 milestone과 자동 완료 장벽도 구별한다.

같은 실제 담당·일정·의존구조의 계획은 제목이나 임시 ID가 달라도 중복 제거한다. 공정성 점수는 해석과 토론의 참고이며, 생성 단계에서 최소 분산안을 자동 승자로 선택하지 않는다.

현재 기본 설정은 raw 조합 500개, 유효 후보 최대 50개, shortlist 최대 10개다. 후보 부족은 수학적 불가능성의 증명이 아니다. \texttt{search\_exhausted}는 해당 구조·탐색 예산 안에서 유효 후보를 확보하지 못했다는 뜻이다.

\newpage
\section{Discussion 확대: 어떤 LLM이 무엇을 판단하는가}
\begin{center}
\begin{tikzpicture}[node distance=5mm]
\node[box] (a) {\textbf{개인 초기 평가 / Assessment}\\본인 카드 + 같은 검증 후보 → 후보별 의사·공정성·카드 적합성 점수};
\node[box,below=of a] (b) {\textbf{코드: 토론 shortlist 구성}\\참가자별 추천 슬롯을 순회해 보장 / private 평가는 공개 대화와 분리};
\node[box,below=of b] (c) {\textbf{개인 발언 / Speech}\\발화자 → 수신자 → 후보 ID → 질문·답변·선호·거절·공개 조건};
\node[box,below=of c] (d) {\textbf{다음 개인 발언}\\이전 발화·거절 이유·조건을 반영 / 라운드마다 순서 회전 / 필요 시 참여자 초대};
\node[box,fill=green,below=of d] (e) {\textbf{새 후보가 필요한 공개 조건인가?}\\조건 + requires\_new\_candidate → 제한된 보충 생성 → Hive 검증·재평가};
\node[box,below=of e] (f) {\textbf{최종 평가 / Assessment → 당사자 Consent}\\최종 점수와 거절 여부 확인 / 정확한 계획 버전마다 새 동의};
\foreach \a/\b in {a/b,b/c,c/d,d/e,e/f} \draw[arr] (\a)--(\b);
\draw[arr] (d.east)--++(7mm,0)|- node[pos=.25,right,font=\small]{다음 턴} (c.east);
\end{tikzpicture}
\end{center}
\textbf{그림 2.} 한 라운드는 한 번의 LLM 호출이 아니다. 초기·최종 평가는 후보를 최대 8개씩 나눠 호출하고, 공개 토론은 참여자별 발언으로 진행한다. 후보별 최종 동의 호출은 평가와 별도다.

\textbf{정보 경계.} 개인 호출에는 자신의 카드와 역할 MD만 주입한다. 다른 사람의 카드 원문·비공개 초기 평가 사유는 공개 토론 입력과 일반 run 본문에 넣지 않는다. 업무 현황, 공개 발언, 공개된 조건은 공유할 수 있다. 프롬프트와 카드 버전을 run 시작 시 고정한다.

\textbf{현재 추천 집계 방식.} 영향 당사자의 최종 평가에 REJECT가 없는 후보를 우선 남기고, 전체 참가자 점수의 평균에 해당하는 합계로 추천 순서를 정한다. 후보별 당사자 Consent에서도 거절이 없어야 최종 후보에 포함된다. 이 집계 규칙은 현재 구현의 구체 방식이며, 사람 공정성의 정답이나 자동 적용을 의미하지 않는다.

\textbf{재동의.} 수정된 계획은 새로운 plan hash와 검증 대상이 된다. 이전 후보의 수락을 새 계획에 자동 이식하지 않는다. 일정·범위·영향 당사자가 바뀌면 재검증과 필요한 재동의를 거친다.

\newpage
\section{PM 선택, 본인 확인, 안전한 적용}
정상적인 추천 완료는 \texttt{awaiting\_pm\_selection}이다. 후보가 하나여도 PM이 명시적으로 선택해야 한다. PM은 전체 반려와 수정 요청도 할 수 있다. 후보가 0개인 실패 전달 \texttt{waiting\_pm}과 구분한다.

PM 선택에는 후보 ID, plan hash와 보드 snapshot을 결합한다. 권한 밖 위임 또는 명시 확인 대상은 \texttt{NEEDS\_INPUT}으로 본인 확인을 기다린다. PM 선택이 개인의 거절을 수락으로 바꾸지 않는다. 확인 대기 중 보드나 카드가 바뀌면 기존 선택을 그대로 적용하지 않고 다시 검사한다. 적용은 트랜잭션 안에서 수행하고 실제 저장 결과도 Hive로 확인한다.

\begin{tabularx}{\linewidth}{@{}p{43mm}Y@{}}\toprule
상태·종료 사유 & 해석\\\midrule
\texttt{running} & 후보 생성·평가·토론·동의 진행 중\\
\texttt{awaiting\_pm\_selection} & 정상 추천 완료, PM 최종 선택 대기\\
\texttt{waiting\_confirm} & 선택안에 필요한 당사자 사람 확인 대기\\
\texttt{applied} & 필요한 동의와 최신 검증 후 실제 보드 반영\\
\texttt{search\_exhausted} & 제한된 탐색에서 후보를 확보하지 못함\\
\texttt{no\_consented\_candidate} & 최종 당사자 비거절 후보가 남지 않음\\
모델·호출·시간 한도 종료 & 모델 오류와 후보 부재를 구분하고 PM에게 요약 전달\\\bottomrule
\end{tabularx}

기본 discussion은 2라운드이며 설정 범위는 1--15다. run 전체에는 기본 120회 실제 호출 예산과 600초 한도를 둔다. 전송 재시도도 예산에 포함한다. 생성 응답 형식 오류는 구체적인 필드 오류를 주고 한 번 복구한다. 평가 후보 ID가 잘못되면 기대 ID·누락·잘못된 ID를 명시해 한 번 다시 요청하며 재실패 시 중단한다. 반복 오류를 무제한 재시도하지 않는다.

\section{화면, API와 기록}
제공된 \texttt{front/Plan Bee Front.html}과 어댑터가 FastAPI에 연결된다. 화면에서 PM / Dev / Art / GD / QA 역할을 바꿔 체험한다. 예시 이슈 버튼은 입력창만 채우며, 보내기를 눌러야 실행한다. 역할 전환 자체는 LLM을 호출하지 않는다.

\begin{tabularx}{\linewidth}{@{}p{69mm}Y@{}}\toprule
API & 목적\\\midrule
\texttt{POST /candidate-runs} & 업무와 줄글 이슈로 run 시작\\
\texttt{POST /runs/\{id\}/advance} & 저장된 상태에서 다음 단계 진행\\
\texttt{POST /runs/\{id\}/pm-selection} & PM 선택·반려·수정 요청\\
\texttt{POST /confirmations/\{id\}} & 해당 본인의 확인 또는 거절\\
\texttt{GET /state}, \texttt{GET /board} & 현재 보드와 제약 결과\\\bottomrule
\end{tabularx}

실행 기록은 \texttt{runs/v2/live/}와 \texttt{runs/v2/mock/}에 날짜·시간·고유값 파일명으로 분리한다. 초기·턴별·최종 milestone, 라운드, 발화자·수신자, 후보와 사유, Hive 검사, 동의 상태·종료 사유를 추적한다. 공개 기록과 비공개 카드·평가 저장을 구분한다. 화면의 역할 헤더는 로컬 체험용 제어이며 실제 로그인 인증 체계를 대신하지 않는다.

\newpage
\section{사람 순위와의 별도 비교 평가}
운영 PM 선택과 블라인드 사람 순위 평가는 분리한다. 고정된 동일 후보 집합에서 공정성과 본인 카드 적합성을 별도 축으로 비교한다. 사람과 LLM에게 제공되는 카드·공개 정보 조건을 맞추고, 사람에게 에이전트 선호를 먼저 노출하지 않는다.

후보 $i$의 사람 순위를 $r_h(i)$, 에이전트 순위를 $r_a(i)$라 할 때 rank-softmax로 분포를 만든다. 동률은 평균 순위로 처리한다.
\[
p_h(i)=\frac{\exp(-(r_h(i)-1)/\tau)}{\sum_j\exp(-(r_h(j)-1)/\tau)},\qquad
q_a(i)=\frac{\exp(-(r_a(i)-1)/\tau)}{\sum_j\exp(-(r_a(j)-1)/\tau)}.
\]
\[
\operatorname{CE}(p_h,q_a)=-\sum_i p_h(i)\log q_a(i).
\]
기본 $\tau=1$이며 낮은 CE가 사람 분포에 더 가깝다. 사람 entropy와 KL, 후보 수, 관측 coverage, 균등 baseline $\log K$, Kendall $\tau_b$ 등을 함께 해석한다. CE는 순위 전체의 순열 likelihood가 아니다. 미평가 후보를 임의로 최하위 처리하지 않으며, 후보가 하나면 순위 비교를 수행하지 않는다. 실제 사람 라벨 없이 점수를 만들어 보고하지 않는다.

\section{현재 연결된 다섯 시나리오}
\begin{tabularx}{\linewidth}{@{}p{13mm}p{30mm}Y@{}}\toprule
구분 & 핵심 이슈 & 검토할 계획·승인\\\midrule
T1 & 일정 이동 & 6시간 Polish를 다음 월요일로 이동, 소요량 불변·Dev 중심 검토\\
T2 & 업무 재분배 & UI 로직은 Dev, 에셋 연결은 Art. Art 신규 인수에 본인 확인\\
T3 & 범위 이관 & 선택 기능과 연결된 잔여 FX를 Beta로 이관, 필수 Core·QA 유지\\
T4 & 복합 조정 & 임시 Integration / 최종 교체와 입력 의존성 재구성. 선택 Retest 생략 시 QA 확인\\
T5 & 현재 제약 충돌 & 필수 수정 뒤 QA 시간을 확보할 후보가 없으면 Dev·QA 근거와 PM 검토 전달\\\bottomrule
\end{tabularx}
원문은 \texttt{seed/m4\_t*/source.md}, 초기 보드·Task·카드는 \texttt{team.yaml}, 이슈는 \texttt{feed.yaml}에 둔다. 원문의 시간 불일치를 정규화한 부분은 시나리오 문서에 공개한다. T4는 Retest를 유지하면서 기한을 충족하는 유효 대안도 있으므로 생략을 강제하지 않는다.

\section{구현 근거와 현재 제한}
주요 구현은 \texttt{app/candidates/engine.py}의 상태 머신, \texttt{generator.py}의 후보 확장·중복 제거, \texttt{plans.py}의 preview·적용, \texttt{app/hive.py}의 일정·제약 검사다. 개인 호출 어댑터는 \texttt{app/llm.py}, 역할별 MD는 \texttt{prompts/agents/}, 생성·토론 MD는 \texttt{prompts/candidate\_generation/}, \texttt{prompts/discussion/}에 있다.

중단 전 v006 전체 회귀 테스트는 199개 통과했고, v007의 관련 테스트 15개가 통과했다. v007 전체 테스트는 사용자 중단 요청 시 170개 통과 상태에서 중단되어 전체 완료로 집계하지 않는다. Live에서도 구조화 응답 오류와 불필요한 본인 확인 해석이 남아 있다. v007 부분 실행은 T1·T2가 PM 대기, T3가 생성 형식 오류, T4가 토론 중 사용자 요청으로 중단되었다. 따라서 모든 Live 시나리오가 안정적으로 성공했다고 주장하지 않는다.

후보 의미의 완전성, 제한 탐색의 해 공간 커버리지, LLM의 카드 해석 일관성은 여전히 한계다. 과거 프롬프트와 실제 실험은 \texttt{past\_versions/agent\_prompts/}에 보존한다. PM 대기 도달, 사람 승인 완료, 실제 적용 성공, 사람 선호 일치도는 서로 다른 지표로 보고해야 한다.
\end{document}
